\subsection{Central Observations}
Our findings highlight three central observations regarding ensemble robustness under distribution shift.

First, ensemble improvements depend on diversity among individual model errors \citep{dietterichEnsembleMethodsMachine2000}.
Uniform averaging provides limited gains when model errors are highly correlated.
The oracle analysis demonstrates that robustness benefits depend on the presence of disagreement under shift.
When experts succeed on different subsets of samples, ensemble potential increases; when they fail together, aggregation offers little advantage.

Second, architectural diversity is a primary driver of such complementary failure modes.
Under real-world shift, heterogeneous ensembles exhibit larger oracle gaps than homogeneous CNN ensembles.
This suggests that distinct inductive biases degrade differently under distribution changes.
As a result, heterogeneous ensembles maintain robustness in regimes where homogeneous ensembles experience correlated collapse.

Third, our mixture-of-experts experiments reveal a limitation of naive routing strategies.
Although routing improves performance relative to uniform averaging, the learned gating mechanism concentrates the majority of its weight on a single dominant expert.
The breakdown of routing decisions shows that only a limited fraction of errors correspond to oracle cases, while a substantial portion arises from hard failures where no expert is correct.

Taken together, these results suggest that architectural diversity alone is insufficient to guarantee effective specialization.
In the absence of explicit incentives, routing mechanisms tend to prioritize globally strong experts rather than conditionally coordinating complementary ones.
Consequently, ensemble robustness under domain shift is constrained both by the diversity of available experts and by the ability of the aggregating mechanisms to exploit that diversity.

Future work may investigate routing objectives that explicitly encourage expert specialization or incorporate input-level uncertainty cues.
Alternatively, increasing expert diversity through data-centric or representation-level strategies may reduce the proportion of hard failures observed under real-world shift.